\writeSectionFrame{\presentationTitle}{UML-Klassendiagramm}
\note{
    Das Klassendiagramm ist das wesentliche Diagramm für Entwurfsmuster und mit diesem werden wir uns am ausführlichsten beschäftigen.
}

\begin{frame}{Das Klassendiagramm}
	\begin{block}{beantwortet die Frage}
		Wie sind die Daten und das Verhalten meines Systems strukturiert?
	\end{block}
\end{frame}

\begin{frame}{Klasse}
    \begin{center}
        \begin{tikzpicture}[scale=0.8, transform shape]
            \umlclass[x=6, y=0]{Klassenname}{
                - attribut1: typ\\
                - attribut2: typ\\
            }{
                + methodeA(parameter1: typ, parameter2:typ) \\
                + methodeB():rueckgabetyp \\
            }
        \end{tikzpicture}
    \end{center}
\end{frame}
\note{
    Klassen werden als Kasten notiert. Der Klassenname steht ganz oben, dann folgen die Attribute und anschließend die Methoden. Die Datentypen werden anders aufgeschrieben als in Java: Sie werden hinter dem Variablenname bzw. dem Methodennamen aufgeschrieben, getrennt durch einen Doppelpunkt. 
}

\begin{frame}{Zugriffsmodifikatoren}
	\begin{itemize}
		\item + : public
		\item -: private
		\item \#: protected
	\end{itemize}
\end{frame}

\begin{frame}{Beispiel}
\begin{tikzpicture}[scale=0.8, transform shape]
    \umlclass[x=6, y=0]{Flugzeug}{
        - pilot: String\\
        - hoehe: double\\
        - geschwindigkeit: double\\
    }{
        + starten() \\
        + landen() \\
        + beschleunigen(menge: double): double\\
        + verlangsamen(menge: double): double\\
        + aufsteigen(menge: double): double\\
        + absteigen(menge: double): double\\
        + pilotWechseln(pilot: String)\\
    }
\end{tikzpicture}
\end{frame}

\begin{frame}{Beziehungen zwischen Klassen}
	\begin{itemize}
		\item Ungerichtete Assozation
        \item Gerichtete Assozation
        \item Vererbung
		\item Aggregation
		\item Komposition
	\end{itemize}
\end{frame}

\begin{frame}{Ungerichtete Assoziation}
    \begin{center}
        \begin{tikzpicture}[scale=0.8, transform shape]
                \umlclass[x=6, y=0]{Klasse1}{
                }{
                }
            
                \umlclass[x=10, y=0]{Klasse2}{
                }{
                }
            
                \umlassoc{Klasse1}{Klasse2}
        \end{tikzpicture}
    \end{center}
\end{frame}
\note{
    Klassen können Abhängigkeiten untereinander aufweisen. Diese werden mit einem durchgezogenen Strich zwischen den beiden beteiligten Klassen notiert. An den Seiten können sog. Multiplizitäten stehen, ähnlich wie bei ER-Diagrammen, z.B. 0, 0 .. 3, 1 .. * Eine ungerichtete Assoziation kann in beide Richtungen gelesen werden.
}

\begin{frame}{Beispiel}
    \begin{center}
        \begin{tikzpicture}[scale=0.8, transform shape]
            \umlclass[x=6, y=0]{Flugzeug}{
            }{
            }

            \umlclass[x=12, y=0]{Pilot}{
            }{
            }

            \umlassoc[mult1=*, mult2=1..2]{Flugzeug}{Pilot}
        \end{tikzpicture}        
    \end{center}
\end{frame}

\begin{frame}{Gerichtete Assoziation}
    \begin{center}
        \begin{tikzpicture}[scale=0.8, transform shape]
                \umlclass[x=6, y=0]{Klasse1}{
                }{
                }
            
                \umlclass[x=10, y=0]{Klasse2}{
                }{
                }
            
                \umluniassoc{Klasse1}{Klasse2}
        \end{tikzpicture}
    \end{center}
\end{frame}
\note{
    Eine gerichtete Assoziation kann nur in die Richtung des Pfeils gelesen werden. Multiplizitäten können auch hier angegeben werden.
}

\begin{frame}{Beispiel}
    \begin{center}
        \begin{tikzpicture}[scale=0.8, transform shape]
            \umlclass[x=6, y=0]{Flugzeug}{
            }{
            }

            \umlclass[x=12, y=0]{Pilot}{
            }{
            }

            \umluniassoc[mult1=*, mult2=1..2]{Flugzeug}{Pilot}
        \end{tikzpicture}        
    \end{center}
\end{frame}

\begin{frame}{Vererbung}
    \begin{center}
        \begin{tikzpicture}[scale=0.8, transform shape]
                \umlclass[x=6, y=0]{Klasse1}{
                }{
                }
            
                \umlclass[x=10, y=0]{Klasse2}{
                }{
                }
            
                \umlinherit{Klasse1}{Klasse2}
        \end{tikzpicture}
    \end{center}
\end{frame}
\note{
    Eine Vererbung wird dargestellt als Pfeil mit geschlossener leerer Spitze. Die Pfeilspitze zeigt in Richtung der Basisklasse.
}

\begin{frame}{Beispiel}
    \begin{center}
        \begin{tikzpicture}[scale=0.8, transform shape]
            \umlclass[x=6, y=0]{Flugzeug}{
            }{
            }

            \umlclass[x=6, y=-3]{Frachtflugzeug}{
            }{
            }

            \umlclass[x=12, y=0]{Pilot}{
            }{
            }

            \umluniassoc[mult1=*, mult2=1..2]{Flugzeug}{Pilot}
            \umlinherit{Frachtflugzeug}{Flugzeug}
        \end{tikzpicture}        
    \end{center}
\end{frame}

\begin{frame}{Implementierung}
    \begin{center}
        \begin{tikzpicture}[scale=0.8, transform shape]
                \umlclass[type=interface, x=6, y=0]{Klasse1}{
                }{
                }
            
                \umlclass[x=10, y=0]{Klasse2}{
                }{
                }
            
                \umlimpl{Klasse1}{Klasse2}
        \end{tikzpicture}
    \end{center}
\end{frame}
\note{
    Eine Implementierung wird durch einen Pfeil mit gestrichelter Linie symbolisiert und ein Interface wird entsprechend markiert. Da dies nicht in allen Programmiersprachen möglich ist, wird dieser Zusammenhang auch oft einfach nur als Vererbung dargestellt, wenn das Diagramm von der verwendeten Sprache unabhängig sein soll. Ob sie dann als Vererbung oder als Implementierung realisiert wird, ist dann ein Detail der tatsächlichen Umsetzung. Bezieht sich das Diagramm tatsächlich auf ein Java-Programm, kann implements ohne weiteres verwendet werden.
}

\begin{frame}{Beispiel}
    \begin{center}
        \begin{tikzpicture}[scale=0.8, transform shape]
            \umlclass[type=interface,x=6, y=3]{Fahrzeug}{
            }{
            }
        
            \umlclass[x=6, y=0]{Flugzeug}{
            }{
            }

            \umlclass[x=6, y=-3]{Frachtflugzeug}{
            }{
            }

            \umlclass[x=12, y=0]{Pilot}{
            }{
            }

            \umluniassoc[mult1=*, mult2=1..2]{Flugzeug}{Pilot}
            \umlinherit{Frachtflugzeug}{Flugzeug}
            \umlimpl{Flugzeug}{Fahrzeug}
        \end{tikzpicture}        
    \end{center}
\end{frame}

\begin{frame}{Aggregation}
    \begin{center}
        \begin{tikzpicture}[scale=0.8, transform shape]
                \umlclass[x=6, y=0]{Klasse1}{
                }{
                }
            
                \umlclass[x=10, y=0]{Klasse2}{
                }{
                }
            
                \umlaggreg{Klasse1}{Klasse2}
        \end{tikzpicture}
    \end{center}
\end{frame}
\note{
    Eine Aggregation ist eine Teil-Ganzes-Beziehung, d.h. damit wird ausgedrückt, dass Klasse2 ein Teil von Klasse1 ist. Wichtig ist dabei, dass die beiden Klassen unabhängig voneinander existieren können. 
}

\begin{frame}{Beispiel}
    \begin{center}
        \begin{tikzpicture}[scale=0.8, transform shape]
            \umlclass[type=interface,x=6, y=3]{Fahrzeug}{
            }{
            }

            \umlclass[x=6, y=0]{Flugzeug}{
            }{
            }
        
            \umlclass[x=3, y=0]{Flughafen}{
            }{
            }

            \umlclass[x=6, y=-3]{Frachtflugzeug}{
            }{
            }

            \umlclass[x=12, y=0]{Pilot}{
            }{
            }

            \umluniassoc[mult1=*, mult2=1..2]{Flugzeug}{Pilot}
            \umlinherit{Frachtflugzeug}{Flugzeug}
            \umlimpl{Flugzeug}{Fahrzeug}
            \umlaggreg{Flughafen}{Flugzeug}
        \end{tikzpicture}        
    \end{center}
\end{frame}
\note{
    In diesem Fall enthält ein Flughafen Flugzeuge, aber die Flugzeuge können auch existieren, wenn es den Flughafen nicht mehr gibt. 
}

\begin{frame}{Komposition}
    \begin{center}
        \begin{tikzpicture}[scale=0.8, transform shape]
                \umlclass[x=6, y=0]{Klasse1}{
                }{
                }
            
                \umlclass[x=10, y=0]{Klasse2}{
                }{
                }
            
                \umlcompo{Klasse1}{Klasse2}
        \end{tikzpicture}
    \end{center}
\end{frame}
\note{
    Eine Komposition ist eine verstärkte Aggregation, d.h. die Teile können ohne das Ganze nicht existieren.
}

\begin{frame}{Beispiel}
    \begin{center}
        \begin{tikzpicture}[scale=0.8, transform shape]
            \umlclass[type=interface,x=6, y=3]{Fahrzeug}{
            }{
            }

            \umlclass[x=3, y=0]{Flughafen}{
            }{
            }

            \umlclass[x=3, y=-3]{Terminal}{
            }{
            }
        
            \umlclass[x=6, y=0]{Flugzeug}{
            }{
            }

            \umlclass[x=6, y=-3]{Frachtflugzeug}{
            }{
            }

            \umlclass[x=12, y=0]{Pilot}{
            }{
            }

            \umluniassoc[mult1=*, mult2=1..2]{Flugzeug}{Pilot}
            \umlinherit{Frachtflugzeug}{Flugzeug}
            \umlimpl{Flugzeug}{Fahrzeug}
            \umlaggreg{Flughafen}{Flugzeug}
            \umlcompo{Flughafen}{Terminal}
        \end{tikzpicture}        
    \end{center}
\end{frame}
\note{
    Wenn der Flughafen nicht mehr existiert, ergeben die Terminals keinen Sinn mehr und können dann ebenfalls entfernt werden.
}

\begin{frame}{Abhängigkeit}
    \begin{center}
        \begin{tikzpicture}[scale=0.8, transform shape]
                \umlclass[x=6, y=0]{Klasse1}{
                }{
                }
            
                \umlclass[x=10, y=0]{Klasse2}{
                }{
                }
            
                \umldep{Klasse1}{Klasse2}
        \end{tikzpicture}
    \end{center}
\end{frame}
\note{
    Eine Abhängigkeit ist etwas schwammmiger als die anderen Beziehungen: eine gerichtete Beziehung zwischen zwei Modellenlementen, bei der ein Element auf das andere angewiesen ist. Dies bedeutet, dass eine Änderung am Ziel (dem abhängigen Element) eine Auswirkung auf das Quell-Element (das abhängige Element) haben könnte. Wobei auch hier gilt: Abhängigkeiten gibt es in einem Programm sehr viele, aber nicht alle sind wichtig im Klassendiagramm. Das Klassendiagramm betont diejenigen, die wirklich eine hohe Bedeutung haben. Die Abhängigkeit zeigt mit einem gestrichelten Pfeil vom abhängigen Element auf das Element, von dem es abhängt. Eine Abhängigkeit ist immer gerichtet. Abhängigkeiten sind schwache Beziehungen im Vergleich zu anderen Beziehungen wie Assoziationen oder Vererbungen. Das Quell-Element kennt das Ziel-Element und ist auf dieses angewiesen, aber nicht umgekehrt. In UML wird diese Art von Abhängigkeit verwendet, um lose Kopplungen anzugeben, bei denen eine Klasse auf eine andere angewiesen ist, ohne eine direkte, dauerhafte Beziehung zu haben (wie bei Assoziationen oder Aggregationen).
}

\begin{frame}{Beispiel}
    \begin{center}
        \begin{tikzpicture}[scale=0.8, transform shape]
            \umlclass[type=interface,x=6, y=3]{Fahrzeug}{
            }{
            }

            \umlclass[x=3, y=0]{Flughafen}{
            }{
            }

            \umlclass[x=0, y=-3]{Terminal}{
            }{
            }
        
            \umlclass[x=6, y=0]{Flugzeug}{
            }{
            }

            \umlclass[x=10, y=-3]{Frachtflugzeug}{
            }{
            }

            \umlclass[x=5, y=-3]{Passagierflugzeug}{
            }{
                +boardAt(terminal: Terminal)
            }

            \umlclass[x=12, y=0]{Pilot}{
            }{
            }

            \umluniassoc[mult1=*, mult2=1..2]{Flugzeug}{Pilot}
            \umlinherit{Frachtflugzeug}{Flugzeug}
            \umlinherit{Passagierflugzeug}{Flugzeug}
            \umlimpl{Flugzeug}{Fahrzeug}
            \umlaggreg{Flughafen}{Flugzeug}
            \umlcompo{Flughafen}{Terminal}
            \umldep{Passagierflugzeug}{Terminal}
        \end{tikzpicture}        
    \end{center}
\end{frame}
\note{
    Ein Passagierflugzeug ist eine Spezialisierung der allgemeinen Klasse Flugzeug. Es landet und startet am Flughafen und verwendet dabei ein bestimmtes Terminal für das Boarding und Deboarding von Passagieren. Dabei ist das Flugzeug auf das Terminal angewiesen, da es bestimmte Terminal-Ressourcen (z.B. Gates) nutzen muss, aber das Terminal selbst kennt das Flugzeug nicht direkt. Eine Abhängigkeit ist hier nützlich, um diese lose Beziehung auszudrücken. Das Flugzeug kennt das Terminal, weil es darauf angewiesen ist, bestimmte Dienste zu nutzen (z.B. Boarding-Gate, Check-in-Services). Änderungen an der Klasse Terminal könnten Auswirkungen auf die Klasse Flugzeug haben (z.B. Änderung der Methode zur Zuordnung eines Gates). Das Flugzeug hat möglicherweise eine Methode wie plane.boardAt(Terminal terminal), die angibt, dass es an einem bestimmten Terminal Passagiere boarden möchte.
}

\begin{frame}{Unterschied Komposition und Aggregation}
	\begin{block}{Aggregation}
        \begin{itemize}
            \item Eine Teil-Ganzes-Beziehung, Teile können unabhängig vom Ganzen existieren.
            \item Wird nur verwendet, wenn die Besteht-aus-Beziehung betont werden soll
            \item Realisierung z.B. mit einer List oder einem Array möglich.
        \end{itemize}
	\end{block}

    \begin{block}{Komposition}
        \begin{itemize}
            \item Eine Teil-Ganzes-Beziehung, Teile können NICHT unabhängig vom Ganzen existieren.
            \item Betont, dass bei der Implementierung besonderer Wert auf die Zerstörung der Teile gelegt werden muss
            \item Eher Ausnahmefall, die Regel ist die Aggregation
        \end{itemize}
	\end{block}
\end{frame}

\begin{frame}{Semantische Unterschiede zwischen den Beziehungen}
	\begin{block}{Vererbung und Implementierung}
        Ist-Ein-Beziehung
	\end{block}
    \begin{block}{Aggregation und Komposition}
        Teil-Ganzes-Beziehung
	\end{block}

    \begin{block}{Assoziation}
        Beziehungsart nicht näher bestimmt bzw. in dieser Betrachtung nicht wichtig
	\end{block}

    \begin{block}{Abhängigkeit}
        Lose und nicht dauerhafte Beziehung (z.B. durch einfachen Methodenaufruf)
	\end{block} 
\end{frame}

\begin{frame}{Granularität von UML-Klassen-Diagrammen}
	\begin{exampleblock}{Wichtiger Hinweis}
        \begin{itemize}
            \item Stellen Sie nicht alles im Klassendiagramm dar
            \item Ein Klassendiagramm enthält nur die Informationen, die für eine bestimmte Betrachtung notwendig sind
            \item Schreiben Sie nie alle Attribute und Methoden
            \item Klassendiagramme vereinfachen den Sachverhalt
            \item Auch Datentypen dürfen weggelassen werden, wenn unwichtig
        \end{itemize}
	\end{exampleblock}
\end{frame}

\begin{frame}{Verschiedene Sichten auf das Projekt}
	\begin{itemize}
		\item Auftraggeber: wenig Details, High-Level
        \item Entwickler: viele Details, sehr technisch
        \item Datenbankentwickler: detailliert in den Datenbankklassen, vereinfacht in den anderen Schichten
        \item Projektleiter: detaillierter aber dennoch grober Überblick
	\end{itemize}
\end{frame}
\note{
    Ein Softwareprojekt hat verschiedene Zielgruppen, die Informationen über das Projekt benötigen. Aber nicht jede Zielgruppe benötigt alle Informationen. Deswegen werden die Klassendiagramme entsprechend der Zielgruppe vereinfacht oder detaillierter. Dem Auftraggeber reicht ein sehr grober Überblick über das Projekt. Er muss die Attribute und Methoden nicht kennen, nicht einmal alle Klasse. Hier könnte man bspw. die grafische Oberfläche mit einer einzigen Klasse symbolisieren. Ein Entwickler muss viele Details kennen, aber vielleicht auch nicht alle. Ein Datenbankentwickler muss nicht den genauen Aufbau der grafischen Oberfläche kennen oder der Logikschicht. Der Projektleiter braucht genauere Informationen als der Auftraggeber, aber auch ihm reicht im Grunde ein grober Überblick. 
}

\begin{frame}{Flexible Sichten auf das Projekt}
	\begin{itemize}
		\item Wenn Informationen fehlen: hinzufügen
        \item Wenn zu detailliert: Informationen entfernen
        \item Es gibt nicht ein Klassendiagramm pro Projekt
        \item Es gibt eine 1:n-Beziehung zwischen Projekt und Klassendiagrammen
	\end{itemize}
\end{frame}
\note{
    In der Flexibilität liegt der große Vorteil von UML-Diagrammen. Was zu viel ist für eine Zielgruppe, kann man flexibel weglassen. Genauso schnell hat man Informationen eingefügt, wenn doch noch Nachfragen kommen. Aus diesem Grund gibt es zu einem Projekt oder sogar zu einem Sachverhalt mehrere Diagramme aus verschiedenen Blickwinkeln und für verschiedene Zielgruppen.
}

\begin{frame}{Kann man Informationen einfach weglassen?}
    \begin{alertblock}{Problem}
        Woher weiß ich, welche Informationen ich ohne Probleme weglassen kann?
    \end{alertblock}
	\begin{exampleblock}{Antwort}
	    Bedürfnisse der Zielgruppe verstehen lernen und natürlich durch Erfahrung
	\end{exampleblock}
\end{frame}
\note{
    Gerade am Anfang ist es nicht einfach zu verstehen, welche Informationen man weglassen darf. Da hilft dann nur ein Verständnis für die entsprechende Zielgruppe und natürlich Erfahrung sowohl in der Programmierung als auch in Planung und Konzeptionierung von Software
}

\begin{frame}{Einsatzmöglichkeiten von Klassendiagrammen}
	\begin{itemize}
		\item Dokumentation
        \item Planung
        \item Analyse mit dem Auftraggeber
		\item Unterstützung der Entwicklung
		\item Kommunikation im Team und mit anderen Teams
	\end{itemize}
\end{frame}

\begin{frame}{Einsatzmöglichkeiten von Klassendiagrammen}
	\begin{exampleblock}{Komplexe Zusammenhänge ...}
	    ... sind grafisch leichter darstellbar, verständlich und diskutierbar als im Quellcode
	\end{exampleblock}
\end{frame}

\begin{frame}{Klassendiagramm bei Entwurfsmustern}
        \begin{itemize}
            \item Entwurfsmuster beschreiben nur einen Lösungsansatz
            \item Klassendiagramme enthalten nur unbedingt notwendige Information für diese Lösung
            \item Weggelassen werden:
                \begin{itemize}
                    \item meistens Attribute
                    \item unnötige Methoden
                    \item unnötige Klassen
                    \item meistens Datentypen
                \end{itemize}
            \item Wichtig sind die Beziehungen zwischen Klassen
        \end{itemize}
\end{frame}

\begin{frame}{Entwurfsmuster als Schablonen}
	\begin{block}{Klassendiagramme in Entwurfsmustern}
        sind sehr allgemeingültige Vorlagen, die auf konkrete Probleme angepasst werden müssen. Deswegen werden sie in der Implementierung oft:
        \begin{enumerate}
            \item erweitert um weitere Elemente
            \item modifiziert an die Fähigkeiten einer Programmiersprache (z.B. Interface statt Vererbung oder Verwendung von Enums)
        \end{enumerate}
        Manchmal werden Klassen im Klassendiagramm auch auf mehrere Implementierungsklassen aufgeteilt
	\end{block}
\end{frame}